% ============================================================================
% Fichier  : partie4/P04_C14_forensique_reseau.tex
% Titre    : Forensique Réseau
% Auteur   : MINKA MI NGUIDJOÏ Thierry Emmanuel
% ============================================================================
% Sources : RFC 3227 ; NIST SP 800-86 ; MITRE ATT&CK. Opération FARO.
% Le bureau d'analyse n'a pas d'outil de capture graphique : on travaille sur
% des journaux de flux exportés (CSV).
% ============================================================================

\chapter{Forensique Réseau}
\label{chap:reseau}
\chaptermark{Forensique réseau}

\epigraph{%
    « Les fichiers mentent parfois~; les flux, rarement. »%
}{Principe du cours}

\niveauChapitre{Intermédiaire à expert}{Chapitre~\ref{chap:memoire}~; notions
de réseaux (TCP/IP).}

\begin{boxobjectifs}
\begin{itemize}
    \item Identifier les sources de traces réseau (flux, journaux, DNS).
    \item Repérer une exfiltration et un canal de pilotage.
    \item Lire des enregistrements de flux exportés.
    \item Relier les traces réseau à la chronologie système.
    \item Rédiger un constat réseau.
\end{itemize}
\end{boxobjectifs}

\bigskip

% ============================================================================
\section*{Question directrice : où sont partis les 40 Go, et par quel canal l'attaquant pilotait-il~?}
% ============================================================================

L'exfiltration et le pilotage laissent des traces réseau~: connexions
sortantes inhabituelles, requêtes DNS répétées, volumes anormaux. Ce chapitre
apprend à les lire, à partir des journaux de flux exportés (le bureau
d'analyse n'a pas d'outil de capture graphique).

\medskip
\textbf{Plan du chapitre.}
\begin{enumerate}
    \item Les sources de traces réseau (section~\ref{sec:c14-sources}).
    \item Exfiltration et pilotage (section~\ref{sec:c14-exfil}).
    \item Lire des flux exportés (section~\ref{sec:c14-flux}).
    \item Recouper avec le système (section~\ref{sec:c14-recouper}).
\end{enumerate}

\separateur

% ============================================================================
\section{Les Sources de Traces Réseau}
\label{sec:c14-sources}
% ============================================================================

\begin{tableaucours}{L{3.6cm}L{8.6cm}}{%
    \tableauHeader{Source} &
    \tableauHeader{Ce qu'elle révèle}
}
Journaux de flux (pare-feu) & Qui parle à qui, quand, combien \\
Journaux du mandataire web & Destinations web, volumes, agents \\
Journaux DNS & Noms résolus, requêtes répétées ou suspectes \\
Journaux d'authentification distante & Accès VPN, bureau à distance \\
\end{tableaucours}
\vspace{-0.8em}
\begin{center}{\small\itshape Tableau~14.1. Sources de traces réseau}\end{center}

\begin{boxalert}
\textbf{Contrainte du cours.} Le bureau d'analyse de NDAP SUKULU ne dispose
pas d'outil de capture graphique.

\smallskip
\textbf{Règle~:} on travaille sur des \emph{journaux de flux exportés} (au
format tableur), ce qui suffit à établir les faits et reste parfaitement
opposable. L'essentiel n'est pas l'outil, mais la lecture méthodique.
\end{boxalert}

% ============================================================================
\section{Exfiltration et Pilotage}
\label{sec:c14-exfil}
% ============================================================================

Deux signatures réseau trahissent une intrusion active~: l'\termecle{exfiltration}
(un volume sortant anormal vers une destination inhabituelle) et le
\termecle{canal de pilotage} (de petites connexions régulières vers un serveur
distant, qui commandent le code malveillant). Les deux se repèrent dans les
journaux de flux.

% ============================================================================
\section{Lire des Flux Exportés}
\label{sec:c14-flux}
% ============================================================================

\begin{boxprocedure}
\textbf{Analyser des journaux de flux}
\begin{enumerate}
    \item Agréger par couple source-destination et par volume.
    \item Repérer les destinations externes inhabituelles.
    \item Repérer les connexions régulières de faible volume (pilotage).
    \item Relever les pics de volume sortant (exfiltration).
    \item Recaler les horodatages sur la chronologie système.
\end{enumerate}
\end{boxprocedure}

\begin{boxscenario}{Opération FARO. Les 40 Go de la nuit du 4 mars}
\textbf{Constatation.} Les journaux de flux montrent, dans la nuit du 4~mars,
un transfert sortant de 40~Go depuis \texttt{FS-01} vers une adresse externe,
précédé, depuis le 24~février, de petites connexions régulières (toutes les
10~minutes) vers la même adresse.

\medskip
\textbf{Analyse} (confiance élevée). Les petites connexions régulières
constituent le canal de pilotage~; le transfert de 40~Go, l'exfiltration. Les
dates recoupent la chronologie système (création du compte le 24, accès aux
partages le 4). Le faisceau est cohérent.
\end{boxscenario}

% ============================================================================
\section{Recouper avec le Système}
\label{sec:c14-recouper}
% ============================================================================

Une trace réseau isolée prouve peu~; recoupée avec le disque et la mémoire,
elle devient un faisceau solide. La force de l'enquête tient à la convergence
des trois scènes (système, mémoire, réseau) sur une même chronologie.

\begin{boxobjet}{Flux d'exfiltration du 4 mars}
\textbf{Nature~:} enregistrement de flux (journal du pare-feu, exporté).
\textbf{Source~:} copie du journal (condensat au dossier). \textbf{Pertinence~:}
transfert de 40~Go de \texttt{FS-01} vers une adresse externe.
\textbf{Fiabilité~:} intègre, horodatage recalé. \textbf{Opposabilité~:}
chaîne de custody continue, recoupé avec la chronologie système.
\end{boxobjet}

\begin{boxmetier}{Analyste forensique et analyste SOC}
\textbf{Faire converger les scènes.} L'analyste ne se contente pas d'une
source~: il recoupe réseau, système et mémoire. Une conclusion tenue par une
seule trace est fragile~; tenue par trois convergentes, elle est opposable.
\end{boxmetier}

\separateur

\begin{boxresume}
\textbf{Ce qu'il faut retenir de ce chapitre~:}
\begin{itemize}
    \item \textbf{Flux, mandataire, DNS}~: les sources réseau.
    \item \textbf{Exfiltration} (gros volume sortant) et \textbf{pilotage}
          (petites connexions régulières).
    \item \textbf{Des journaux de flux exportés} suffisent à établir les
          faits.
    \item \textbf{On recoupe} réseau, système et mémoire.
    \item \textbf{La convergence} fait la force de la preuve.
\end{itemize}
\end{boxresume}

\begin{boxinfo}
\textbf{Points de contrôle de l'élève.} Je sais~: citer les sources réseau~;
reconnaître exfiltration et pilotage~; lire des flux exportés~; recouper avec
la chronologie~; rédiger un constat réseau.
\end{boxinfo}

\bigskip

\begin{boxexo}[title={\ding{45}\quad Exercice~14.1. Sources \niveaubase}]
Quelle source révèle le volume d'un transfert sortant~? Les noms résolus~?
\end{boxexo}

\begin{boxexo}[title={\ding{45}\quad Exercice~14.2. Exfiltration \niveaubase}]
Quelle signature réseau trahit une exfiltration~?
\end{boxexo}

\begin{boxexo}[title={\ding{45}\quad Exercice~14.3. Pilotage \niveaubase}]
Comment se repère un canal de pilotage dans les flux~?
\end{boxexo}

\begin{boxexo}[title={\ding{45}\quad Exercice~14.4. Sans capture \niveaumoyen}]
Pourquoi des journaux de flux exportés suffisent-ils, sans capture graphique~?
\end{boxexo}

\begin{boxexo}[title={\ding{45}\quad Exercice~14.5. FARO \niveaumoyen}]
Dans FARO, distinguez le canal de pilotage de l'exfiltration.
\end{boxexo}

\begin{boxexo}[title={\ding{45}\quad Exercice~14.6. Recouper \niveaumoyen}]
Comment les flux du 4~mars recoupent-ils la chronologie système~?
\end{boxexo}

\begin{boxexo}[title={\ding{45}\quad Exercice~14.7. Faisceau \niveaumoyen}]
Pourquoi une trace réseau isolée prouve-t-elle peu~?
\end{boxexo}

\begin{boxexo}[title={\ding{45}\quad Exercice~14.8. Constat réseau \niveaumoyen}]
Rédigez le constat du flux d'exfiltration du 4~mars.
\end{boxexo}

\begin{boxexo}[title={\ding{45}\quad Exercice~14.9. Analyse de flux \niveauexpert}]
À partir d'un journal de flux exporté, proposez la démarche pour isoler
l'exfiltration.
\end{boxexo}

\begin{boxexo}[title={\ding{45}\quad Exercice~14.10. Convergence \niveauexpert}]
Construisez le faisceau de preuves de FARO en faisant converger les trois
scènes.
\end{boxexo}

\bigskip

\begin{boxinfo}
\textbf{Pour aller plus loin.}
\begin{itemize}
    \item RFC~3227~\cite{rfc3227}~; NIST SP~800-86~\cite{nist80086}~;
          MITRE ATT\&CK~\cite{mitreattack}.
\end{itemize}
\end{boxinfo}

\begin{boxinfo}
\textbf{Transition.} L'enquête établie, il faut agir~: la Partie~V traite de
la réponse à incident, à commencer par son cycle
(Chapitre~\ref{chap:cycle}).
\end{boxinfo}
